% part: intuitionistic-logic
% chapter: introduction
% section: syntax

\documentclass[../../../include/open-logic-section]{subfiles}

\begin{document}

\olfileid{int}{int}{syn}

\olsection{అంతఃప్రజ్ఞావాద తర్కపు వాక్యనిర్మాణం}

అంతఃప్రజ్ఞావాద తర్కపు వాక్యనిర్మాణం ప్రతిజ్ఞావాక్య
తర్కానిదే. సాంప్రదాయిక ప్రతిజ్ఞావాక్య తర్కంలో
కొన్ని సంధాయకాలను ఇతర సంధాయకాలతో నిర్వచించవచ్చు.
ఉదాహరణకు, $!A \lif !B$ను $\lnot !A \lor !B$తో,
లేదా $!A \lor !B$ను $\lnot(\lnot !A \land \lnot
!B)$తో నిర్వచించవచ్చు. అందుకే సాంప్రదాయిక తర్కపు
వివరణల్లో కొన్ని సంధాయకాలను ఆ నిర్వచనాల
సంక్షిప్త రూపాలుగా ప్రవేశపెడతారు. అంతఃప్రజ్ఞావాద
తర్కంలో మాత్రం రెండు మినహాయింపులు తప్ప అలా కాదు:
$\lnot !A$ను $!A \lif \lfalse$కు సంక్షిప్త రూపంగా
నిర్వచించవచ్చు; తరచూ అలానే నిర్వచిస్తారు. అలా
అయితే $\lfalse$ను మాత్రం నిర్వచిత సంకేతంగా
తీసుకోకూడదు!{} అలాగే $!A \liff !B$ను సాంప్రదాయిక
తర్కంలోలాగే $(!A \lif !B) \land (!B \lif
!A)$గా నిర్వచించవచ్చు.

ప్రతిజ్ఞావాక్య అంతఃప్రజ్ఞావాద తర్కపు
\tetoken{సూత్రాలను}{!!^{formula}s},
\emph{\tetoken{ప్రతిజ్ఞావాక్య చరాలు}{!!{propositional variable}s}},
ప్రతిజ్ఞావాక్య స్థిరాంకం~$\lfalse$ నుంచి
\emph{తార్కిక సంధాయకాల} సహాయంతో నిర్మిస్తాం. ఇవి:

\begin{enumerate}
\item \tetoken{ప్రతిజ్ఞావాక్య చరాల}{!!{propositional variable}s}
  \tetoken{ఒక అనంతంగా లెక్కించదగిన}{!!^a{denumerable}}
  సమితి~$\PVar$: $\Obj p_0$, $\Obj p_1$, \dots
  \item \tetoken{అసత్యం}{!!{falsity}} కోసం ప్రతిజ్ఞావాక్య స్థిరాంకం~$\lfalse$.
\item తార్కిక సంధాయకాలు: $\land$
  (సంయోజనం), $\lor$ (విడియోజనం), $\lif$ (సోపాధికం)
\item విరామ సంకేతాలు: (, ), అలాగే కామా.
\end{enumerate}

\begin{defn}[సూత్రం]
\ollabel{defn:formulas} ప్రతిజ్ఞావాక్య అంతఃప్రజ్ఞావాద
తర్కపు \emph{\tetoken{సూత్రాల}{!!{formula}s}} సమితి~$\Frm[L_0]$ను
ఇలా ఆగమన పద్ధతిలో నిర్వచిస్తాం:
\begin{enumerate}
\item $\lfalse$ ఒక అణు \tetoken{సూత్రం}{!!{formula}}.

\item ప్రతి \tetoken{ప్రతిజ్ఞావాక్య చరం}{!!{propositional variable}}~$\Obj p_i$
  ఒక అణు \tetoken{సూత్రం}{!!{formula}}.

\item $!A$, $!B$ \tetoken{సూత్రాలు}{!!{formula}s} అయితే,
  $(!A \land !B)$ \tetoken{ఒక సూత్రం}{!!a{formula}}.

\item $!A$, $!B$ \tetoken{సూత్రాలు}{!!{formula}s} అయితే,
  $(!A \lor !B)$ \tetoken{ఒక సూత్రం}{!!a{formula}}.

\item $!A$, $!B$ \tetoken{సూత్రాలు}{!!{formula}s} అయితే,
  $(!A \lif !B)$ \tetoken{ఒక సూత్రం}{!!a{formula}}.

\tagitem{limitClause}{ఇవి తప్ప మరేదీ \tetoken{సూత్రం}{!!a{formula}} కాదు.}{}
\end{enumerate}
\end{defn}

పైన ప్రవేశపెట్టిన ప్రాథమిక సంధాయకాలతో పాటు,
కింది \emph{నిర్వచిత} సంకేతాలను కూడా వాడతాం:
$\lnot$ (నిరాకరణ), $\liff$
(\tetoken{ద్విసోపాధికం}{!!{biconditional}}).
నిర్వచిత కారకాలతో నిర్మించిన సూత్రాలను ఇలా
అర్థం చేసుకోవాలి:
\begin{enumerate}
\item $\lnot !A$ అనేది $!A \lif \lfalse$కు సంక్షిప్త రూపం.

\item $!A \liff !B$ అనేది $(!A \lif !B) \land (!B
  \lif !A)$కు సంక్షిప్త రూపం.
\end{enumerate}

$\lnot$ను అధికారికంగా సంక్షిప్త రూపంగా తీసుకున్నా,
కొన్నిసార్లు అది ప్రాథమిక సంకేతమైనట్లుగా
$\lnot$కు నిర్వచనాల్లో స్పష్టమైన నియమాలు, నిబంధనలు ఇస్తాం.
ముఖ్యంగా అభ్యాస సమస్యలను చెప్పడానికి ఇది ఉపయోగపడుతుంది.

\end{document}
